\scanpage{411}\mspage{150}
tout au moins fidèle, de sorte que ça prouvera
que $T^2$ induit l'identité dans $\mathrm{Hot}$ également.

Je reviens aux questions $1^\circ$ à $4^\circ$ d'il y a qques
jours. Pour $1^\circ$), il faut prouver que
\[
p\colon\underline{\mathrm{Ch}}^{\deleted{\unsure{W}}}_{\infty}(X)\longrightarrow X\times X
\]
est lisse et propre et une $\mathrm{HOT}$-fibration. Par
dualité, \inserted{$p$} propre résultera de $p$ lisse. Mais
$\underline{\mathrm{Ch}}^{\deleted{W}}_{\infty}(X)$ est $\varinjlim$ des $\underline{\mathrm{Ch}}^{\deleted{W}}_{n,m}(X)$, avec $n\longrightarrow-\infty$, $m\longrightarrow+\infty$,
et où on peut se borner à $n$ et $m$ \uline{pairs}. Or
on voit que pour $n,m$ pairs, \unsure{i.e.|avec} $\tau_{n,m}$ à flèche
initiale \deleted{positive}, directe, à flèche finale rétrograde,
$\underline{\mathrm{Ch}}_{n,m}^{\deleted{W}}(X)\simeq\underline{\mathrm{Ch}}_{\tau_{n,m}}^{\deleted{W}}(X)$ est \inserted{Cat-}\uline{fibrée} sur $X\times X$,
donc \inserted{\deleted{\unsure{une}}}\deleted{$\mathrm{HOT}$-fibration\unsure{,}} lisse (i.e. $\mathrm{HOT}$ lisse),
\unsure{ceci|cela} \unsure{passe} aux $\varinjlim$.
% slanted
\margin{\uline{NB} $\underline{\mathrm{Ch}}_{\infty}(X)$ n'est ni fibré ni cofibré sur $X\times X$ -- car les foncteurs d'inclusion entre les $\underline{\mathrm{Ch}}_{n,m}$ ne sont pas cartésiens \unsure{resp.} \inserted{p.ex.} cocartésiens.}
Ainsi $\underline{\mathrm{Ch}}^{\deleted{W}}_{\infty}(X)\xrightarrow{p}X\times X$ est à la fois propre et
lisse.$^{*}$
% slanted
\margin{$^{*}$12.11. Mais ça \unsure{implique|signifie|infirme} déjà que c'est une fibration !}
Dire que c'est une $\mathrm{HOT}$-fibration, revient
à dire que pour
\[
(a,b)\xrightarrow{(u,v)}(a',b')
\]
\unsure{dans} $X\times X$, la flèche correspondante
\[
\underline{\mathrm{Ch}}^{\deleted{W}}_{\infty}(X)_{/(a,b)}\longrightarrow
\underline{\mathrm{Ch}}^{\deleted{W}}_{\infty}(X)_{/(a',b')}
\]
est une $\mathrm{Hot}$-équivalence. Or, comme $p$ est
propre, les inclusions des fibres en $(a,b)$ et en
$(a',b')$ dans l'un et l'autre membre sont des
$\mathrm{Hot}$ équivalences :
% archive sheet number: 410
